\chapter{Review of existing literature on maritime surveillance systems, covering ship detection, multi-object tracking, and trajectory prediction methods}

\section{Introduction}
Maritime surveillance has evolved from traditional radar and manual monitoring systems to advanced, AI-driven solutions. The increasing complexity of maritime activities, the rise in illegal operations, and the need for environmental protection have driven research towards more robust, automated, and intelligent surveillance systems. This chapter presents a comprehensive review of the literature, focusing on ship detection, multi-object tracking, and trajectory prediction, and highlights the challenges, gaps, and future directions in the field.


This fast development of deep learning and computer vision has changed the way maritime surveillance research is conducted, and a series of studies on intelligent, automated solutions to traditional radar and manual surveillance systems have emerged. The chapter reviews the literature on the three main technical aspects of the proposed system, namely, ship detection, multi-object tracking, and trajectory prediction, from the latest publications in the most prestigious journals such as Expert Systems with Applications, Engineering Applications of Artificial Intelligence, Journal of Marine Science and Engineering and Scientific Reports. The characteristics of each reviewed work are analyzed for their major contributions and methodological decisions as well as their limitations, and the works are specifically analyzed in terms of how well they address the complexities of the real world in maritime surveillance. The gaps as a whole that have been identified from literature are the technical and operational justification for the end-to-end integrated system proposed in this project.

\section{Details about the Literature Review of the papers conducted for this project}

\subsection{Ship Detection in Maritime Environments}
Ship detection is the foundational step in maritime surveillance. Early approaches relied on classical image processing and radar, but these methods struggled with clutter, occlusion, and varying weather conditions. The introduction of deep learning, especially Convolutional Neural Networks (CNNs), revolutionized detection accuracy. State-of-the-art detectors such as YOLOv5, YOLOv7, and YOLOv8 have been adapted for maritime use. Li and Wang (2025) proposed EGM-YOLOv8, which reduced model parameters by 13.57\% while improving recall, making it suitable for real-time applications on resource-constrained platforms. Jian et al. (2025) demonstrated that large-scale, domain-specific training data significantly enhances detection in complex sea scenes. However, most detection models are limited to single-frame analysis and do not address temporal consistency or identity preservation.

Recent works have also explored transformer-based architectures and hybrid models, which show promise in handling scale variation and adverse weather. However, these models often require significant computational resources and are not always suitable for real-time deployment on edge devices. The literature also highlights the importance of robust annotation and diverse datasets, as models trained on limited data often fail to generalize to real-world scenarios.

\subsection{Multi-Object Tracking (MOT) in Maritime Surveillance}
Tracking multiple vessels across video frames is challenging due to occlusion, re-entry, scale changes, and dynamic backgrounds. Traditional tracking-by-detection approaches, such as SORT and DeepSORT, have been widely used. Ma et al. (2026) extended DeepSORT with OCR for hull number recognition, improving identity consistency but requiring high-resolution imagery. Hao et al. (2025) combined YOLOv7tiny with StrongSORT and trajectory association, reducing identity switches and enabling real-time tracking on embedded hardware.

Recent advances include appearance-based Re-Identification (ReID) using deep features, motion modeling with Kalman filters, and sensor fusion with AIS data. However, most systems are evaluated in controlled environments and struggle with real-world challenges such as heavy occlusion, low visibility, and crowded scenes. The lack of standardized benchmarks and annotated datasets for maritime MOT remains a significant barrier.

\subsection{Trajectory Prediction}
Predicting vessel trajectories is crucial for proactive surveillance, anomaly detection, and collision avoidance. Early methods relied on linear models and Kalman filters, which are efficient but limited to linear motion. Recent research leverages Recurrent Neural Networks (RNNs), especially Long Short-Term Memory (LSTM) networks, to capture complex temporal dependencies. These models can predict future positions based on historical tracks, enabling early warning of anomalous or illegal behavior.

Despite progress, trajectory prediction in maritime environments faces challenges such as irregular sampling, missing data, and non-linear vessel maneuvers. Integrating contextual information (e.g., weather, traffic lanes, and regulations) remains an open research area. Few works address the fusion of video-based and AIS-based predictions, which is essential for comprehensive situational awareness.

\subsection{Sensor Fusion and System Integration}
Modern maritime surveillance systems increasingly integrate multiple data sources, including video, AIS, and radar. Sensor fusion techniques, such as Kalman and particle filters, are used to combine high-frequency but noisy video data with sparse but accurate AIS observations. This integration improves robustness and accuracy but introduces challenges in data association, synchronization, and real-time processing.

End-to-end pipelines that combine detection, tracking, prediction, and visualization are rare in the literature. Most works focus on individual components, and few address the operational requirements of real-time, large-scale deployment. The proposed system in this project aims to fill this gap by integrating all modules into a unified, scalable pipeline.

\subsection{Challenges and Research Gaps}
The literature reveals several persistent challenges:
\begin{itemize}
	\item \textbf{Occlusion and Clutter:} Vessels often overlap or are partially hidden, making detection and tracking difficult.
	\item \textbf{Scale Variation:} Ships appear at different scales due to camera perspective and distance.
	\item \textbf{Adverse Weather:} Fog, rain, and low light degrade image quality and model performance.
	\item \textbf{Data Scarcity:} Lack of large, annotated maritime datasets limits model generalization.
	\item \textbf{Real-Time Constraints:} Many advanced models are computationally intensive and unsuitable for real-time use.
	\item \textbf{Integration:} Few systems offer seamless integration of detection, tracking, and prediction in a single pipeline.
\end{itemize}

\subsection{Future Directions}
Future research should focus on:
\begin{itemize}
	\item Developing lightweight, robust models for deployment on edge devices.
	\item Creating large, diverse, and publicly available maritime datasets.
	\item Improving sensor fusion techniques for better situational awareness.
	\item Incorporating contextual and semantic information into prediction models.
	\item Standardizing benchmarks and evaluation protocols for fair comparison.
\end{itemize}

The current studies on maritime surveillance based on deep learning show a fragmented and evolving literature, with each study tackling a particular aspect of the overall surveillance tasks without providing a coherent end-to-end approach. Although the proposed multi-object ship tracking system is lightweight and competitive to other tracking systems in real-time, it is still limited to ship detection performance due to using an improved YOLOv7tiny backbone, and it shows poor robustness in extreme occlusion while has no predictability at all. Although highly useful as a reference, Zhang et al. suggest no novel system, and do not even touch on the integration of detection, tracking, and prediction, which remains the main hurdle in the application of deep learning to real-world maritime surveillance, and the key challenges of the maritime domain: occlusion, extreme scale variation, water surface clutter, and adverse weather degradation. 
Although the model parameters were reduced by 13.57\% compared to YOLOv8, EGM-YOLOv8 was still lightweight, showing excellent real-time detection performance in maritime scenes, especially on small and distant targets, it only performed detection tasks and lacked temporal awareness, identity tracking, and trajectory analysis in the subsequent video frames. While Jian et al. have contributed to the performance gains in complex sea scenes with high detection accuracy, and generalisation across a variety of maritime scenes and states, the system's restrictions on computation and single-frame detection (with no tracking, trajectory extraction, or predictive module) severely limit its applicability to the most basic stage of the surveillance pipeline. 
To address the issue, Ma et al. developed a framework called DeepSORT which is extended by adding an OCR module to extract hull identification number as an additional hull feature to be used for re-identification, but a high resolution image is needed for the system to work, the hull number needs to be visible, and it cannot be used for long-range surveillance or in adverse weather conditions, and there is no trajectory prediction function. The system is designed for low visibility inland waterway conditions (such as fog, backlight and water reflections) and is based on dehazing-aware feature extraction and illumination-robust normalisation, making it robust even in foggy scenarios but limited to inland waterway detection scenarios, without any tracking, identity management, motion analysis or path prediction. 
A common and critical deficiency identified in all reviewed systems is the lack of a real-time, unified pipeline that can perform ship detection, persistent multi-object tracking with identity-preserving, motion parameter estimation and future trajectory prediction — a deficiency which fuels and sets the boundaries of the system proposed in this project.

\section{Programming Language used}

The suggested maritime surveillance system is developed in full using Python, a high-level interpreted language which now is the de facto standard for Artificial Intelligence, Deep Learning and Computer Vision research and development. The simplicity and eloquence of Python's syntax are also driving its use in these applications — by streamlining the process of coding mathematically intensive algorithms, engineers can concentrate on designing their systems and can model the logic of the systems without having to worry about all the implementation details. However, Python is also the native interface language of almost all current, widely-used deep learning and computer vision libraries, providing seamless interoperability between all components of the proposed pipeline. It is also the most practical and productive option for an academic project this size and interdisciplinary in scope because of its extensive community support, rich documentation and its ability to be rapidly prototyped.

\section{Tools and Technologies used}

The system is based on a properly selected set of open-source libraries and frameworks that play specific roles in the pipeline. PyTorch is the backbone of deep learning, supporting dynamic building of computational graphs, automatic differentiation and native support of the CUDA libraries for the training and inference of models on GPUs. Building on PyTorch, the Ultralytics YOLOv8 library offers an easy to use API to create ship detection models, set up training and deployment parameters, and use the model for real-time inferences. Video ingestion is done with OpenCV, a computer vision library that has been highly optimised for real-time frame-by-frame processing pipelines, needed for live surveillance applications, such as video ingestion.OpenCV is a highly optimised computer vision library, which supports frame-by-frame processing pipelines, and therefore is indispensable for live surveillance applications, including video ingestion. 
The DeepOCSORT tracking module is a fully integrated tracking module that fuses directly with the detection output to preserve vessel identities from frame to frame via appearance-based Re-Identification and Kalman filter state estimation. Using NumPy and Pandas, numerical computation, array manipulation, and calculations of motion parameters (centroids) are carried out, which is fast and vectorised data handling of large trajectory datasets. Trajectory and detection visualisations are created in the video output using Matplotlib, geospatial vessel position mapping is created using Folium and interactive maps of tracked vessel positions are generated. All of the pipeline is speeded up via NVIDIA CUDA, allowing the system to process the detection, tracking, and prediction tasks on real time. Testing and experimentations are done in Jupyter Notebook and Visual Studio Code, enabling iterative testing and integration of modules before rollout in the complete pipeline.


\section{Summary}

This chapter had conducted a comprehensive survey of the existing literature in maritime surveillance, systematically reviewing the published literature across the marine detection field, multi-object tracking and trajectory prediction fields from top-notch academic journals. In summary, the summarized works in this review show that deep-learning and computer vision approaches have rendered great achievements in the development of robust, lightweight detection systems, as well as appearance-based tracking models and techniques for maritime monitoring purposes which are resilience and adaptable to different adverse environmental conditions. But there is a hardly autarkic mistake which prevails throughout all above-mentioned publications: up to now there is no satisfactory system combining ship detection with keeping ships identity within a track record of callsigns or unique identifiers, motion analysis into a track, and forecasting the trajectory of the sought ships in a single real-time pipeline. The programming tools and language used in the development of this project were identified and justified as the most suitable and effective tools available to meet this gap, these tools were: Python, PyTorch, YOLOv8, DeepOCSORT, LSTM, OpenCV and Folium. The identified research gaps and the technological alternatives proposed in this chapter in respect to the maritime surveillance system, directly feed into the discussion of the overall system design and method proposed for their implementation in the subsequent chapter, which more broadly outlines the required architectural choices and design rationale behind this end-to-end maritime surveillance system.
